\documentclass[11pt]{article}
\usepackage[T1]{fontenc}
\usepackage{lmodern,amsmath,amssymb,amsthm,geometry,booktabs,longtable,microtype}
\geometry{margin=1in}
\usepackage[hidelinks]{hyperref}
\hypersetup{pdftitle={Uniform real-rootedness of the (132,213)-avoiding permutation-polytope family},pdfauthor={mxym}}
\newtheorem{theorem}{Theorem}[section]
\newtheorem{lemma}[theorem]{Lemma}
\newtheorem{corollary}[theorem]{Corollary}
\newtheorem{proposition}[theorem]{Proposition}
\theoremstyle{remark}\newtheorem{remark}[theorem]{Remark}
\newcommand{\Prb}{\mathbb P}
\newcommand{\E}{\mathbb E}
\newcommand{\Var}{\operatorname{Var}}
\newcommand{\Rpoly}{\mathcal R}
\newcommand{\Hcal}{\mathcal H}
\newcommand{\Bcal}{\mathcal B}
\newcommand{\Fcal}{\mathcal F}
\newcommand{\EE}{\mathcal E}
\newcommand{\norm}[1]{\left\|#1\right\|_x}
\newcommand{\dd}{\,dt}
\newcommand{\cc}{2-\sqrt3}
\title{Uniform real-rootedness of the $(132,213)$-avoiding\\permutation-polytope family}
\author{mxym}
\date{8 October 2026}
\begin{document}
\maketitle
\begin{abstract}
We prove that the Ehrhart $h^*$-polynomial of the convex hull of the
$(132,213)$-avoiding permutations in $S_d$ has only negative real zeros for
every $d\ge3$. A uniform analytic argument proves positivity of the repaired
Eulerian tails for every $d\ge1001$. Combining this with de Castro's finite
Theorem 1.3 for $d\le1000$ resolves Conjecture 10.1 of
arXiv:2609.06096v1. The new argument uses centered Bernoulli-sum comparison,
finite centered moments of an infinite correction series, and degree smoothing
for its large components. Exact rational certificates establish scalar
envelopes on 33 closed intervals and the finite bases of infinite inductions.
The original finite dataset is an explicitly cited prior input, not an
independently reconstructed computation. A short corollary applies Zhang's
known Veronese preservation theorem to all positive integral dilates of this
particular family.
\end{abstract}
\tableofcontents

\section{Statement, dependencies, and notation}\label{sec:inputs}
Let $P_d$ be the convex hull of the permutations in $S_d$ avoiding both
$132$ and $213$, in its induced affine lattice, and write
$H_d(t)=h^*(P_d;t)$. The subscripted polynomials $P_n(z)$ below denote
primitive auxiliary products, not the polytope $P_d$.
\begin{theorem}\label{thm:main}
For every integer $d\ge3$, all zeros of $H_d$ are real and strictly negative.
For $d=1,2$, $H_d=1$.
\end{theorem}
The new analytic result in this paper covers \emph{every} $d\ge1001$.
The remaining dimensions use the prior finite Theorem 1.3 of
de Castro~\cite{decastro}. We do not claim a new general preservation theorem
for other pattern classes, poset permutahedra, or arbitrary lattice polytopes.

We use the following precise prior inputs from the fixed version
\cite{decastro}. Its refined Eulerian coordinates $a_{d,j}$ satisfy the exact
transform in its equation (25). Its Corollary 6.2 says that nonnegative
endpoint coordinates together with
\begin{equation}\label{eq:tails}
 S_{d,r}(c)=\sum_{j=r}^{d-2}c^{j-r}a_{d,j}\ge0
 \quad(1\le r\le d-2),\qquad c=\cc,
\end{equation}
imply real-rootedness for $d\ge4$. The endpoint coordinates are nonnegative,
with $a_{d,0}=1$ and positive terminal coordinate for $d\ge3$.
Its Theorem 1.3 supplies $d\le1000$.
The criterion is sometimes called Proposition 6.2 elsewhere in that source;
we refer to its displayed Corollary 6.2.

The primitive and divisor polynomials, and the signed correction series, are
\begin{align}
 P_n(z)&=\frac{1}{n^2(n-1)!}\prod_{q=1}^{n-1}(2n-q+qz),\label{eq:primitive}\\
 J_n(z)&=\sum_{h\mid n}\frac{\varphi(h)}{h^2}(z-1)^{n-n/h}P_{n/h}(z),\\
 J(u,z)&=\sum_{n\ge1}J_n(z)u^n,\qquad
 \Hcal(u,z)=e^{(z-1)J(u,z)}=\sum_{s\ge0}\Hcal_s(z)u^s.
\end{align}
Their exact marked identity is
\begin{equation}\label{eq:marked}
 \Rpoly_d(z):=\sum_{m=0}^{d-1}a_{d,d-1-m}z^m
 =\sum_{s=0}^{d-1}\frac{d-s}{d}J_{d-s}(z)\Hcal_s(z).
\end{equation}
These identities are equations (34)--(37) of \cite{decastro}.
For $x>0$ let
\begin{align}
 f_x(t)&=\frac{xt}{2+(x-1)t},&
 p&=\frac{x}{1+x},&
 \nu&=\int_0^1f_x(t)\dd,\\
 X_{n,x}&=\sum_{q=1}^{n-1}\operatorname{Bernoulli}(f_x(q/n)),&
 \mu_n&=\E X_{n,x},&
 V_n&=\Var X_{n,x}.
\end{align}
All Bernoulli summands are independent. In particular,
\begin{equation}\label{eq:PB}
 [z^j]P_n(z)x^j=P_n(x)\Prb(X_{n,x}=j),\qquad
 \frac{\mu_n}{1+x}\le V_n\le\mu_n.
\end{equation}
The mean $\mu_d(x)$ increases continuously and strictly from $0$ to $d-1$;
hence $1\le m\le d-2$ determines a unique saddle $\mu_d(x)=m$.

Set
\begin{equation}
 \lambda=\exp\left(1+\int_0^1\log(2+(x-1)t)\dd\right),\quad
 R=\lambda^{-1},\quad q=(1+x)R,\quad E=\Hcal(R,x).
\end{equation}
We also use the following product and norm bounds from Lemma 9.1 of
\cite{decastro}, uniformly for $0\le x\le15$:
\begin{align}
 P_n(x)R^n
 &=\frac{e^{\delta_n}}{\sqrt{4\pi(1+x)}\,n^{5/2}},
 \qquad -1/n<\delta_n<0\quad(n\ge2),\label{eq:product}\\
 \lambda&\ge4,\qquad q\le4/5,\qquad
 \norm{J_n-P_n}R^n\le(4/5)^{n/2}\quad(n\ge1).\label{eq:periodic}
\end{align}
Here $\norm F=\sum_j|[z^j]F|x^j$.
Darroch's theorem~\cite{darroch}, in the form reviewed in
\cite[p.~284]{pitman}, identifies an integer mean of a finite Bernoulli sum
with its largest atom. No prior eventual-dimension conclusion is needed below.

\section{The forward repaired tails}\label{sec:forward}
We sharpen only the mass constant in the forward estimate of
\cite[Lemmas 8.1--8.3 and Proposition 8.4]{decastro}.
Write
\[
 L_n(x)=\frac{1}{n^2(n-1)!}\prod_{q=1}^{n-1}(xn+q),
 \qquad \EE(x)=\frac{(1+x)^{1+x}}{x^x}.
\]
For $5/54\le x\le2$, the Robbins and Riemann bounds give
\[
 L_n(x)\EE(x)^{-n}\le\frac{n^{-5/2}}{x\sqrt{2\pi}}
 <\frac92n^{-5/2}\quad(n\ge2).
\]
At $n=1$ the normalized term is $q_F=\EE(x)^{-1}<49/64$.
Since $\sum_{n\ge2}n^{-5/2}<3/8$ and
$\varphi(h)/h^2\le1/2$ for $h\ge2$, the normalized connected mass satisfies
\[
 A<\left(\frac{49}{64}+\frac92\frac38\right)
       \left(1+\frac{q_F}{2(1-q_F)}\right)
 <\frac{157}{64}\frac{79}{30}=\frac{12403}{1920}<\frac{13}{2}.
\]
The connected series is $K(u,x)=\sum_n k_n(x)u^n$, with
$k_n(x)=\sum_{h\mid n}\varphi(h)h^{-2}L_{n/h}(x)$, and
$\Fcal_d(x)=[u^d]e^{K(u,x)}$. Its coefficients are nonnegative.
The prior connected coefficient bound $5000n^{-5/2}$ remains valid.
For nonnegative logarithmic coefficients bounded by $Cn^{-5/2}$ and mass
$A$, a $k$-part composition has a part at least $n/k$; freely summing the
others yields $Ck^{7/2}A^{k-1}n^{-5/2}$.
Summing its exponential weights, using $k^{7/2}\le k^4$, gives
\[
 [u^n]\exp\left(\sum b_ju^j\right)
 \le Ce^A(A^3+6A^2+7A+1)n^{-5/2}.
\]
Consequently the prior $\exp(140)$ envelope can be replaced by $\exp(30)$:
\[
 5000e^A(A^3+6A^2+7A+1)
 <5000\cdot687e^{13/2}<2^{22}e^{13/2}<e^{30}.
\]
Here is the explicit tail comparison to which the improved envelope applies.
Let $y=\sqrt3-1$, $q_C=5/113$, and $x_C=5/54$. The refined transform gives
\[
 \mathcal A_d(z)=\sum_j a_{d,j}z^j
 =(1-z)^{d-1}\Fcal_d(2z/(1-z)).
\]
Its alternative polynomial expansion in the nonnegative coefficients of
$\Fcal_d$ bounds it on $|z|=q_C$ by
$(1+q_C)^{d-1}\Fcal_d(x_C)$.
Put $\mathcal M=c^{-r}y^{d-1}\Fcal_d(y)>0$.
Subtracting the initial coordinates and the terminal coordinate gives
\[
 S_{d,r}(c)=\mathcal M-c^{-r}\sum_{j<r}c^ja_{d,j}
                  -c^{d-1-r}a_{d,d-1}.
\]
Cauchy's coefficient bound and the geometric sum show that the relative
absolute initial-coordinate loss is at most
\[
 \mathcal R=\frac{q_C}{c-q_C}(c/q_C)^r
       ((1+q_C)/y)^{d-1}\frac{\Fcal_d(x_C)}{\Fcal_d(y)}.
\]
The terminal loss is at most
$\mathcal U=(c/y)^{d-1}\Fcal_d(2)/\Fcal_d(y)$, since
$0<a_{d,d-1}\le\Fcal_d(2)$.
The prior primitive lower bound is
$\Fcal_d(y)\ge d^{-5/2}\EE(y)^d/36$.
The exact inequalities of \cite[Lemma 8.3]{decastro} are
\[
 (c/q_C)^{1/4}\frac{1+q_C}{y}\frac{\EE(x_C)}{\EE(y)}<19/20,
 \qquad \frac cy\frac{\EE(2)}{\EE(y)}<77/100 .
\]
For $r\le d/4$, insert the improved upper envelope and use
$q_C(c-q_C)^{-1}y/(1+q_C)<1$ and $y/c<3$.
The two relative losses are therefore bounded by
$36e^{30}(19/20)^d$ and $108e^{30}(77/100)^d$.
All intermediate bounds in that argument hold already for $d\ge2$; its old
dimension threshold arose only from these final errors. For $d\ge1001$,
$e<4$, $(19/20)^{14}<1/2$, and $(77/100)^3<1/2$ give respective bounds
$2^{-5}$ and $2^{-266}$. Their sum is below one.
\begin{proposition}\label{prop:forward}
For $d\ge1001$ and $1\le r\le\lfloor d/4\rfloor$, $S_{d,r}(\cc)>0$.
\end{proposition}

\section{Centered local comparison}\label{sec:centered}
\subsection{Fourier profiles and a modal lower bound}
For a finite Bernoulli parameter vector $\boldsymbol r$, of mean $\mu$ and
variance $V$, define its centered characteristic function and real profile
\[
 \psi(t)=\E e^{it(X-\mu)},\qquad
 F(y)=\frac1{2\pi}\int_{-\pi}^{\pi}\psi(t)e^{-iyt}\dd .
\]
The profile is real by conjugate symmetry. It equals a point probability
only when $y+\mu$ is an integer; no probability interpretation at other
arguments is needed.

For one parameter put $g_r(t)=e^{-irt}(1-r+re^{it})$.
Taylor's integral remainder gives
\[
 |\partial_rg_r(t)|
 \le|e^{it}-1-it|+r|t||e^{it}-1|\le\tfrac32t^2.
\]
If the other summands in each replacement have variance at least $W\ge0$,
their characteristic product is bounded by $e^{-2W\sin^2(t/2)}$.
Pad the shorter vector by zero parameters and let
$D=\sum_i|p_i-q_i|$. Telescoping proves
\begin{equation}\label{eq:centercompare}
 |F_{\boldsymbol p}(y)-F_{\boldsymbol q}(y)|
 \le\tfrac32D C_2(W),\quad
 C_2(W)=\frac1{2\pi}\int_{-\pi}^{\pi}t^2e^{-2W\sin^2(t/2)}\dd .
\end{equation}
Apply replacements with nonnegative variance increments first, then negative
ones. Every intermediate total variance is at least the smaller endpoint
variance. Removing one coordinate costs at most $1/4$, so one may take
$W=\max(0,\min(V_{\boldsymbol p},V_{\boldsymbol q})-1/4)$.
For $W>0$, using $\sin(t/2)\ge|t|/\pi$ yields
\begin{equation}\label{eq:C2general}
 C_2(W)\le\min\left(\frac{\pi^2}{3},
       \frac{\pi^{5/2}}{4\cdot2^{3/2}}W^{-3/2}\right).
\end{equation}

Linear sine terms in $\operatorname{Im}g_r$ cancel, giving
$|\operatorname{Im}g_r(t)|\le r(1-r)|t|^3/6$.
Telescoping a product against its conjugate gives
$|\operatorname{Im}\psi(t)|\le(V/6)|t|^3
 e^{-2(V-1/4)\sin^2(t/2)}$ when $V>1/4$.
Differentiating the finite integral therefore proves
\begin{align}
 |F(\Delta)-F(0)|&\le A_1(V)|\Delta|
                       +\tfrac12C_2(V)\Delta^2,\label{eq:shift}\\
 A_1(V)&\le V\pi^4/30\quad(V\ge0),\\
 A_1(V)&\le\frac{\pi^{9/2}}{16\cdot2^{5/2}}
                   \frac{V}{(V-1/4)^{5/2}}\quad(V>1/4).\label{eq:A1general}
\end{align}
In particular $A_1(V)<4$ for all $V>0$, and
$A_1(V)<4V^{-3/2}$ for $V\ge1$.

\begin{lemma}\label{lem:mode}
Every integer-valued distribution with finite variance $V$ and largest atom
$M$ satisfies $M\ge(12V+1)^{-1/2}$. If a Bernoulli sum has integer mean $m$,
this bounds $\Prb(X=m)$ from below.
\end{lemma}
\begin{proof}
Adding an independent uniform variable on $[-1/2,1/2]$ gives a density bounded
by $M$, with variance $V+1/12$. Among mass-one densities bounded by $M$, the
density equal to $M$ on the interval of length $1/M$ centered at the mean
minimizes the second moment. Indeed, subtract that density and integrate
against $(t-\mu)^2-1/(4M^2)$: both the inside and outside products are
nonnegative, and the constant term integrates to zero. Its variance is
$1/(12M^2)$. Darroch's integer-mean mode theorem gives the last assertion.
\end{proof}

\subsection{Sharp constants for variance at least twenty}\label{sec:sharp}
\begin{lemma}\label{lem:sharp}
If $V\ge20$, then $C_2(V)<(3/5)V^{-3/2}$ and
$A_1(V)<(1/3)V^{-3/2}$.
\end{lemma}
\begin{proof}
For $|t|\le1$, $\sin(t/2)\ge t/2-t^3/48$ and hence
$2\sin^2(t/2)\ge11t^2/24$. For $1\le|t|\le\pi$ it is at least
$529/1152$. The Gaussian second moment is below $(12/25)V^{-3/2}$,
using $\sqrt{2\pi}>12/5$; squaring verifies the remaining rational comparison.
The other portion is at most $\pi^2e^{-529V/1152}$.
After multiplication by $V^{3/2}$ this decreases for $V\ge20$ and is
below $900/8000=9/80$ at twenty. Thus $12/25+9/80<3/5$.

Put $W=V-1/4\ge79V/80$. The imaginary-part bound makes $A_1$ at most $V/6$
times a normalized fourth-moment integral. Split at $1$ and $\pi/2$.
The first portion, enlarged to the full Gaussian line, is below
$(8/5)W^{-5/2}$. On $[1,\pi/2]$, $\sin t\ge\sin1>5/6$, so
\[
 2\sin^2(t/2)\ge529/1152+(5/6)(t-1).
\]
Since $t^4<7$ and $\pi>3$, this portion is below
$(14/(5W))e^{-529W/1152}$. The last portion is below $100e^{-W}$.
It follows that
\[
 A_1(V)<\frac4{15}VW^{-5/2}
 +\frac7{15}\frac VW e^{-529W/1152}
 +\frac{50}{3}Ve^{-W}.
\]
After multiplication by $V^{3/2}$ these are bounded by
\[
 \frac4{15}\frac{26}{25}
 +\frac7{15}\frac{80}{79}\frac{90}{8000}
 +\frac{50}{3}\frac{1800}{10^8}<\frac13.
\]
Here $(80/79)^3<26/25$, the normalized exponential terms decrease,
$529(79/4)/1152>9$, and $79/4>19$.
The positive sums $\sum_{k=0}^{40}9^k/k!>8000$ and
$\sum_{k=0}^{60}19^k/k!>10^8$ are exact rational witnesses.
\end{proof}

\subsection{The primitive arrays and endpoint error}\label{sec:arrays}
For $n=d-s$, retained parameters satisfy $f_x(q/n)\ge f_x(q/d)$.
The means increase in the dimension: shifting $q$ to $q+1$ embeds the old
array beneath the new one. Thus the padded parameter distance is
\begin{equation}\label{eq:D}
 D=\mu_n-\mu_d+2\sum_{q=n}^{d-1}f_x(q/d)\le2sp.
\end{equation}
Put $e_j=j\nu-\mu_j$. The increasing Riemann bounds give $0\le e_j\le p$.
At $\mu_d=m$ the shift needed by the n-law is
\[
 \Delta=m-k-\mu_n=s\nu-k+\eta,\qquad
 \eta=e_n-e_d=\mu_d-\mu_n-s\nu.
\]
The trapezoid formula writes $e_j=p/2-\epsilon_j$ with
\[
 |\epsilon_j|\le\frac{|f_x'(1)-f_x'(0)|}{8j}.
\]
The errors $\epsilon_j$ have the same sign for all $j$, because $f_x$ is
convex for $x<1$, concave for $x>1$, and linear for $x=1$. Consequently
$|\eta|\le |f_x'(1)-f_x'(0)|/(8n)$, without adding two endpoint bounds.
For $x\le1/10$ this is at most $3x/(16n)$; for $x\le15$ it is at most
$15/(16n)$. Concavity also gives $e_j\ge p/2$ when $x\ge1$.

\section{The positive majorant and its centered moments}\label{sec:majorant}
Define the nonnegative series
\[
 \Bcal(u,z)=(1+z)\sum_{n\ge1}u^n
  \sum_{h\mid n}\frac{\varphi(h)}{h^2}(1+z)^{n-n/h}P_{n/h}(z).
\]
Its exponential majorizes the absolute coefficients of $\Hcal$.
At $z=x$, write
\[
 v_a=(1+x)P_a(x)R^a,\quad
 w_{a,h}=v_a\frac{\varphi(h)}{h^2}q^{(h-1)a},\quad
 \beta=\sum_{a,h}w_{a,h},\quad \beta_1=\sum_{a,h}ha\,w_{a,h}.
\]
A logarithmic component has dimension $n=ha$ and degree law
\begin{equation}\label{eq:componentlaw}
 K=X_{a,x}+\operatorname{Binomial}((h-1)a+1,p).
\end{equation}
Set $\ell=(h-1)a$, $t_a=p-e_a$, $r=p-\nu$, and $T_j=\sum w_{a,h}\ell^j$.
Then $0\le t_a\le p$, $r\ge0$, and
\[
 \E(K-n\nu)=t_a+\ell r.
\]
After including the first extra Bernoulli, there are $a$ parameters of total
mean $a\nu+t_a$. Cauchy's inequality bounds their variance by
$a\nu(1-\nu)+t_a(1-2\nu)-t_a^2/a$. Therefore
\begin{equation}\label{eq:jensen}
 \Var K\le n\nu(1-\nu)+p\max(1-2\nu,0)
                      +\ell\{p(1-p)-\nu(1-\nu)\}.
\end{equation}
Let $M_1=\sum w\,\E(K-n\nu)$ and
$M_2=\sum w\,\E(K-n\nu)^2$. If $t_a\le t_0$ and $r\le r_0$, then
\begin{align}
 M_1&\le t_0\beta+r_0T_1,\label{eq:M1}\\
 M_2&\le \nu(1-\nu)\beta_1+p\max(1-2\nu,0)\beta
 +\max(0,p(1-p)-\nu(1-\nu))T_1\notag\\
 &\hspace{6mm}+t_0^2\beta+2t_0r_0T_1+r_0^2T_2.\label{eq:M2}
\end{align}

Write $c_{s,k}=R^sx^k[u^sz^k]\exp\Bcal(u,z)\ge0$ and $h_s=\sum_kc_{s,k}$.
The exponential formula gives
\begin{align}
 \sum_s h_s&=e^\beta,& \sum_s s h_s&=e^\beta\beta_1,\\
 \sum_{s,k}c_{s,k}(k-s\nu)^2&=e^\beta(M_2+M_1^2),\label{eq:compound}\\
 e^{-\beta}\sum_{s,k}c_{s,k}|k-s\nu|
 &\le \min\left(\sqrt{M_2+M_1^2},\,2\nu\beta_1+M_1\right).\label{eq:absolute}
\end{align}
These are infinite-series statements. First restrict to finite sets of
components, for which the independent compound-Poisson calculation is exact.
All coefficients and the squared weights are nonnegative; increasing component
sets increases their sums. The component means are nonnegative, and the
finite bounds below give finite $M_1,M_2,\beta,\beta_1$. Monotone convergence
proves \eqref{eq:compound}. The last bound follows both from Cauchy and
$|k-s\nu|\le k+s\nu$.
The divergent uncentered second \emph{dimension} moment is never used.

Both $q$ and $v_a$ increase with $x$, since
\[
 \frac{d\log q}{dx}=\frac{p-\nu}{x}\ge0,\qquad
 \frac{d\log v_a}{dx}=\frac{p-e_a}{x}\ge0.
\]
Thus $\beta,\beta_1,T_j,b_n,h_s$ increase too, where
$b_n=[u^n]\Bcal(Ru,x)$.
Finally,
\begin{equation}\label{eq:Jmoment}
 \sum_n n\norm{J_n}R^n\le\frac{\beta_1}{1+x},\qquad
 E=e^{(x-1)J(R,x)}>e^{-\beta}.
\end{equation}

\section{Degree smoothing and coefficient-tail induction}\label{sec:smoothing}
\begin{lemma}\label{lem:smoothing}
Let $v_*=p(1-p)$ and
$c(x)=\min(2/3,(1+x)^2/6)$. Each degree law of a fixed composition of
components of total dimension $s$ has variance at least $c(x)s v_*$.
Every coefficient of the normalized positive majorant of dimension $s\ge1$,
even after convolution with a normalized nonnegative polynomial, is at most
\[
 h_s\frac{\sqrt\pi}{2\sqrt{2c(x)s v_*}}.
\]
\end{lemma}
\begin{proof}
The Bernoulli variance integrand is precisely
\[
 g_x(t)=f_x(t)(1-f_x(t))
       =\frac{xt(2-t)}{[2+(x-1)t]^2}.
\]
For $x\ge1$ it is at least $v_*t(2-t)$; for $x\le1$ it is at least
$(1+x)^2v_*t(2-t)/4$. Include the first extra Bernoulli as the endpoint
$t=1$, and use
\[
 \sum_{q=1}^a(q/a)(2-q/a)=\frac{(a+1)(4a-1)}{6a}\ge\frac{2a}{3}.
\]
The remaining $(h-1)a$ Bernoulli variables each have variance $v_*$,
and $c(x)\le1$. Variances add in fixed compositions. Fourier inversion and
$\sin(|t|/2)\ge |t|/\pi$ bound the largest atom of a Bernoulli sum of variance
$V$ by $\sqrt\pi/(2\sqrt{2V})$. A positive mixture has the same bound,
and convolution is a weighted average of translates.
\end{proof}

\begin{lemma}\label{lem:induction}
If $b_n\le Bn^{-5/2}$ for the needed large indices and $h_j\le Cj^{-5/2}$
is proved for the needed smaller indices, then
\begin{equation}\label{eq:induction}
 h_s\le s^{-5/2}\left[
  \frac{6\beta_1}{s}C+Be^\beta\left(1+\frac{9\beta_1}{s}\right)\right].
\end{equation}
\end{lemma}
\begin{proof}
Use $s h_s=\sum_{n=1}^s n b_nh_{s-n}$. For $n\le s/2$, the ratio
$(1-n/s)^{-5/2}$ is at most $2^{5/2}<6$. For $n>s/2$, put $j=s-n$
and use $(1-j/s)^{-3/2}\le1+9j/s$ on $[0,1/2]$.
Sum $\sum h_j=e^\beta$ and $\sum jh_j=e^\beta\beta_1$.
In particular bases $64\le j\le127$ suffice for induction from $s=128$;
no bound for $j<64$ is needed in the first part.
\end{proof}
If $v_a\le Aa^{-5/2}$ and $q\le q_0<1$, the proper divisor contributions give
\begin{equation}\label{eq:b-tail}
 n^{5/2}b_n\le A+3A n^{3/2}q_0^{n/2}.
\end{equation}
Indeed $\varphi(h)/h^2\le1/h=a/n$, $a\le n/2$ for $h\ge2$, and
$\sum_a a^{-3/2}<3$. The right side decreases in the ranges used below.

\section{A common marked-identity decomposition}\label{sec:split}
Fix $d\ge1001$, an integer $1\le m\le d-2$, and its saddle $x\le15$.
Write $\kappa=[z^m]P_d>0$ and $S=\lfloor d/8\rfloor$.
For $n=d-s$, the primitive multiplier is
\begin{equation}
 A_s=\frac nd\lambda^s\frac{P_n(x)}{P_d(x)}
     =(d/n)^{3/2}e^{\delta_n-\delta_d}.
\end{equation}
For $s\le S$, \eqref{eq:product} gives
\begin{equation}\label{eq:A}
 A_s<5/4,\quad |A_s-1|<\{(9/4)s+3\}/d\quad(s\ge1),\quad A_0=1.
\end{equation}
The derivative of $(1-t)^{-3/2}$ is below $9/4$ for $t\le1/8$.
Also $|\delta_n-\delta_d|<1/n$ and
$e^{1/n}\le(1-1/n)^{-1}$; their multiplier contributes less than $3/d$.
For $n\ge d/2$, the coarser bound is $A_s<4$.

For each signed coefficient of $\Hcal_s$, first compare the profiles at zero
by \eqref{eq:centercompare}, then shift the n-profile by
$\Delta=m-k-\mu_n$ using \eqref{eq:shift}. Divide by the integer-mean d-mode
using Lemma~\ref{lem:mode}. Thus the probability error is bounded by
\begin{equation}\label{eq:relative}
 \sqrt{12V_d+1}\left\{\tfrac32D C_2(W)
            +A_1(V_n)|\Delta|+\tfrac12C_2(V_n)\Delta^2\right\}.
\end{equation}
The intermediate n-profile at zero need not itself be a probability.

To compare \eqref{eq:marked} with $E$, use the small-$s$ primitive terms,
the middle primitive terms $S<s\le d/2$, and the full $J_n$ terms $s>d/2$.
Handle $J_n-P_n$ for every $n\ge d/2$ separately.
Finally bound the infinite tail of $E$ after $s=S$ by $\sum_{s>S}h_s$.
These pieces exhaust the marked identity and its comparison series;
integer cutoffs use $s\le\lfloor d/2\rfloor$ where appropriate.

\section{Small saddles}\label{sec:small}
\begin{proposition}\label{prop:small}
If $d\ge1001$, $\mu_d(x)=m\in\{1,\ldots,d-2\}$, and $0<x\le1/10$,
then $[z^m]\Rpoly_d>0$.
\end{proposition}
\subsection{Scalar bounds and their infinite tails}
At $x=1/10$, $q=(11/20)^{20/9}<53/200<4/15$, certified by ninth powers.
For all smaller $x$, \eqref{eq:product} and $\pi>55/18$ imply
$v_a<(3/10)a^{-5/2}$ for $a\ge2$, while $v_1=q<3/10$.
The exact small-saddle script certifies
\begin{equation}\label{eq:smallcaps}
 \beta<2/5,\quad \beta_1<4/5,\quad T_1<3/25,\quad T_2<1/5.
\end{equation}
For completeness its finite mass sums use $a\le64,h\le16$.
For $t=q^a$, the h tails in the mass and dimension moment are respectively
$t^{16}/(17(1-t))$ and $a t^{16}/(1-t)$.
For $a>64$ their total tails are bounded by
$1/2560$ times $1+q^{65}/(2(1-q^{65}))$, and
$(3/40)/(1-q^{65})$.
The direct periodic-moment envelopes are
\[
 T_1\le\frac{q^2}{1-q}
 +\frac9{80}\frac{q^2}{(1-q)(1-q^2)},\quad
 T_2\le\frac{q^2}{(1-q)^2}
 +\frac9{40}\frac{q^2}{(1-q)(1-q^2)^2}.
\]
The first terms handle $a=1$; the others use
$a^{-3/2}<3/8$ and $a^{-1/2}<3/4$ for $a\ge2$.

Equation~\eqref{eq:b-tail} with $A=3/10,q_0=4/15$ is below $2/5$ at
$n=16$ and decreases thereafter. The exact bases $n<16$ verify this same
$b_n$ bound. Exact $h_s^2s^5<1$ for $1\le s\le23$ and
\eqref{eq:induction} at $s=24$, with $B=2/5,C=1,e^\beta<3/2$, give
\[
 h_s<s^{-5/2}\ (s\ge1),\qquad E>2/3,\qquad e^\beta/E<9/4.
\]
Integral bounds yield $(9/25)x\le\nu<(2/5)x$.
For the lower bound, $\nu/x$ decreases in $x$ and at $1/10$ it equals
$-10/9+(200/81)\log(20/11)$; the positive nine-term log series with
$t=9/31$ proves the claim. For the upper bound its limit at zero is
$2\log2-1<2/5$. The integer mean forces $y=dx>5/2$, and $V_d<(2/5)y$.

\subsection{Centered losses for small components}
Equations~\eqref{eq:M1}--\eqref{eq:M2} give
$M_1\le(13/25)x$ and $M_2+M_1^2<(24/25)x$.
One explicit loose bound is $M_2\le(21/25)x+(21/25)x^2$:
its variance's linear coefficient is at most
$(2/5)(4/5)+2/5+(16/25)(3/25)=0.7968$;
the component squared-mean coefficient is at most
$2/5+2(16/25)(3/25)+(16/25)^2(1/5)=0.63552$.
Together with $M_1^2\le(169/625)x^2$, these are below the stated bound at
$x\le1/10$. Equation~\eqref{eq:absolute} also bounds the absolute first
moment by $e^\beta(29/25)x$.

For $s\le S$, $n\ge7d/8$, so $|\eta|<x/(4d)$. Thus the weighted absolute
and squared shift moments are respectively below
$e^\beta(117/100)x$ and $e^\beta(97/100)x$.
The weighted D bound is $2pe^\beta\beta_1<e^\beta(8/5)x$.
The prefactor error from \eqref{eq:A}, divided by $E$, is below $27/(4d)$.

The elementary bound $V_j\ge(j\nu-p)/(1+x)$ gives, for both needed arrays,
\[
 V_n\ge\frac{63}{220}y-\frac1{11},\qquad
 W\ge\frac{63}{220}y-\frac{15}{44}>0.
\]
The bounds of Section~\ref{sec:centered} imply
\begin{align}
 y\sqrt{(24/5)y+1}\,C_2(W)&<50,\\
 y\sqrt{(24/5)y+1}\,C_2(V_n)&<50,\\
 y\sqrt{(24/5)y+1}\,A_1(V_n)&<100.
\end{align}
For $C_2$, split at $y=7/2$. Below it use $C_2<10/3$ and the square-root
bound $17/4$; above it use $W\ge3y/16$, square root below
$(23/10)\sqrt y$, Gaussian constant below $8/5$, and
$(16/3)^{3/2}<25/2$.
For $A_1$, when $y<5$ the global bound $A_1<4$ gives a normalized bound
below 100. When $y\ge5$, $V_n\ge y/4$ and the bound $4V_n^{-3/2}$
gives a bound below 72, hence below 100.
Using \eqref{eq:relative}, the probability error divided by $E$ is below
\[
 \frac1d\frac54\frac94
 \left\{\frac{12}{5}50+\frac{117}{100}100+\frac{97}{200}50\right\}
 <\frac{735}{d}.
\]
Including the prefactor error gives $742/d$.

\subsection{Middle terms, large terms, and periodic terms}
Since $S\ge d/9$, the omitted E tail divided by $E$ is below
$27d^{-3/2}$. For $S<s\le d/2$,
$V_n\ge(9/55)y-1/11$ and the primitive largest-atom/d-mode ratio is below
5, by squaring the Gaussian bound and using $y>5/2,\pi<4$.
With $A_s<4$, the middle contribution is below $540d^{-3/2}$.

For $s>d/2$, Lemma~\ref{lem:smoothing} has $c\ge1/6$ and $v_*>(4/5)x$.
Its atom bound is below $\sqrt{15/(4sx)}$.
Also \eqref{eq:product} gives $P_d(x)R^d>(5/19)d^{-5/2}$, and the inverse
d-mode is below $(23/10)\sqrt y$.
The full $J_n$ contribution is consequently bounded by
\[
 \frac{35}{2}d^2\sum_{n<d/2}n\norm{J_n}R^n(d-n)^{-3}.
\]
Use \eqref{eq:Jmoment}, $\beta_1<4/5$, $E>2/3$, and
$(d-n)^{-3}<8d^{-3}$ to obtain $169/d$ after division by $E$.

All periodic terms with $n\ge d/2$, using \eqref{eq:periodic}, the same
product and mode bounds, and $e^\beta/E<9/4$, total less than
$35d^3e^{-d/20}<1/1000$. The function decreases for $d>60$.
At $1001$, $e>8/3$ and
$35000\cdot1001^3(3/8)^{50}<1$ are an exact witness.
The overall relative error is therefore
\[
 \frac{912}{d}+\frac{567}{d^{3/2}}+\frac1{1000}
 \le\frac{912}{1001}+\frac{567}{1001\cdot31}+\frac1{1000}<1.
\]
Thus $|[z^m]\Rpoly_d/\kappa-E|<E$, proving Proposition~\ref{prop:small}.

\section{Larger saddles and interval certificates}\label{sec:large}
\begin{proposition}\label{prop:large}
If $d\ge1001$, $\mu_d(x)=m$ is an integer, and $1/10\le x\le15$,
then $[z^m]\Rpoly_d>0$.
\end{proposition}
\subsection{Variance and uniform local comparison}
The variance density is
\[
 v(x)=\int_0^1\frac{xt(2-t)}{[2+(x-1)t]^2}\dd.
\]
For $x\ne1$, direct integration gives
\begin{align}
 \nu(x)&=\frac{x}{x-1}-\frac{2x\log((1+x)/2)}{(x-1)^2},\\
 v(x)&=\frac{2x(x+1)\log((1+x)/2)}{(x-1)^3}
                  -\frac{x(3x+1)}{(x+1)(x-1)^2},\label{eq:vclosed}
\end{align}
with $\nu(1)=1/4,v(1)=1/6$.
Both $\nu/x$ and $v/x$ decrease in $x$, while $\nu$ increases.
The integrand $g_x$ is nonnegative, unimodal and at most $1/4$.
Its total variation is at most $1/2-g_x(1)$, with $g_x(0)=0$.
Comparing the right Riemann sum and then subtracting the endpoint proves
\[
 jv-1/2\le V_j\le jv+1/2.
\]
Every certified interval has $v\ge3/100$, so $dv\ge30.03$.
For $s\le S$, this implies
\[
 W\ge(17/20)dv,\quad V_n\ge(6/7)dv,\quad
 \sqrt{12V_d+1}<(7/2)\sqrt{dv}.
\]
In particular $W,V_n>20$.
Lemma~\ref{lem:sharp} yields
\[
 \sqrt{12V_d+1}\,C_2(W),\
 \sqrt{12V_d+1}\,C_2(V_n)<\frac{27}{10dv},
 \quad
 \sqrt{12V_d+1}\,A_1(V_n)<\frac{3}{2dv}.
\]
The numerical comparisons follow by squaring positive rational quantities.
Section~\ref{sec:arrays} gives $|\eta|<1/500$, since
$15/(16n)\le15/(16\cdot7\cdot1001/8)<1/500$.

\subsection{Whole-interval envelopes}\label{sec:envelopes}
Fix a closed interval $[a,b]\subset[1/10,15]$ from Table~\ref{tab:intervals}.
The endpoint helper gives rational lower and upper bounds on $\nu$ and $v$,
and upper bounds $\beta,B_1,T_1,T_2$ at $b$.
Set
\[
 v_{\rm low}=a\,v(b)_{\rm low}/b,\qquad
 v_{\rm high}=\min(1/4,b\,v(a)_{\rm high}/a).
\]
Let $N_{\rm hi},N_{\rm lo}$ be the maximum and minimum of $t(1-t)$ over the
enclosed $\nu$ range, and $P_{\rm hi},P_{\rm lo}$ the analogous bounds over
the exact endpoint p range. These extrema include a possible crossing of
$1/2$. Put $p_a=a/(1+a),p_b=b/(1+b)$ and
\[
 r_0=\min\left(p_b-\nu(a)_{\rm low},
           \frac{p_b(p_a-\nu(a)_{\rm low})}{p_a}\right),\qquad
 t_0=\begin{cases}p_b/2,&a\ge1,\\p_b,&a<1.\end{cases}
\]
The second r bound follows because
$(p-\nu)/p=\int_0^1 2(1-t)/[2+(x-1)t]\dd$ decreases in x.
The improved $t_0$ follows from concavity when $a\ge1$.
Define
\begin{align}
 M&=t_0\beta+r_0T_1,\\
 Q&=N_{\rm hi}B_1+p_b\max(1-2\nu(a)_{\rm low},0)\beta
          +\max(0,P_{\rm hi}-N_{\rm lo})T_1\notag\\
  &\quad+t_0^2\beta+2t_0r_0T_1+r_0^2T_2+M^2,\\
 U&=\min(\sqrt Q,\,2\nu(b)_{\rm high}B_1+M),\qquad \epsilon=1/500.
\end{align}
Equations~\eqref{eq:compound}--\eqref{eq:absolute} bound the required shifted
moments by $e^\beta(U+\epsilon)$ and
$e^\beta(Q+2\epsilon U+\epsilon^2)$.

For a lower bound $L$ on $\log E$, use $L=-\beta$ when $a<1$.
For $a\ge1$, set $\rho=(x-1)/(x+1)$ and note
\[
 \log E=\rho\sum_{j,h}v_j\frac{\varphi(h)}{h^2}
                                  (\rho q)^{(h-1)j}.
\]
All factors are nonnegative and increasing in x. A positive finite sum at
a is therefore a valid lower bound throughout the interval.
Take certified $R_{\exp}\ge e^{\beta-L}$ and $I_{\exp}\ge e^{-L}$.
The small-component probability and prefactor errors, divided by E, are
bounded by $K_{\rm small}/d$ and $K_{\rm pref}/d$, where
\begin{align}
 K_{\rm small}&=\frac54\frac{R_{\exp}}{v_{\rm low}}
 \left\{\frac{81}{10}p_bB_1+\frac32(U+\epsilon)
       +\frac{27}{20}(Q+2\epsilon U+\epsilon^2)\right\},\\
 K_{\rm pref}&=R_{\exp}\left(\frac94B_1+3\beta\right).
\end{align}
In the first term, $(3/2)\cdot(2p_bB_1)\cdot(27/10)=
(81/10)p_bB_1$; thus the full weighted parameter loss is included.

\subsection{Tail bases and the remaining marked terms}
The exact tail script uses the following caps at four endpoints.
By monotonicity each cap applies to every smaller positive x.
Here $A$ bounds $v_a a^{5/2}$, $q_0$ bounds q, B bounds
$b_nn^{5/2}$ for $n\ge64$, and $e_0$ bounds $e^\beta$.
\begin{center}\small
\begin{tabular}{c c c c c c c c}
\toprule
x&A&$q_0$&B&$\beta$&$\beta_1$&$e_0$&C\\\midrule
1&$41/100$&$3/8$&$1/2$&$7/12$&$6/5$&$9/5$&$3/2$\\
3&$4/7$&$1/2$&$3/5$&$81/100$&$7/4$&$9/4$&2\\
7&$81/100$&$16/25$&$5/6$&$9/8$&$8/3$&$25/8$&4\\
15&$8/7$&$3/4$&$4/3$&$3/2$&4&$9/2$&10\\\bottomrule
\end{tabular}
\end{center}
The script checks $A^2>(1+x)5/62$ using $\pi>31/10$, and $q_0<A$ for
$a=1$. It checks \eqref{eq:b-tail} at 64 and its decreasing continuation,
the exact h bases 64..127, and
$6\beta_1C/128+Be_0(1+9\beta_1/128)<C$.
Lemma~\ref{lem:induction} gives $h_s<Cs^{-5/2}$ for every $s\ge64$.

For $S<s\le d/2$, $V_n\ge dv/2-1/2$; the squared Gaussian atom ratio
to the d-mode is below four squared, since $dv\ge30$.
Using $A_s<4$ and $S\ge d/9$, the middle contribution and omitted E tail
together are below $306CI_{\exp}d^{-3/2}$.

For $s>d/2$, use Lemma~\ref{lem:smoothing} and \eqref{eq:Jmoment}.
Set
\[
 c_{\rm low}=\min(2/3,(1+a)^2/6),\qquad
 U_P=\sqrt{4(22/7)(1+b)}\,\frac{1001}{1000}.
\]
Equation~\eqref{eq:product} and $e^{1/d}\le1001/1000$ give the inverse
primitive scale bound $U_Pd^{5/2}$.
Since $\sqrt\pi/(2\sqrt2)<16/25$ and $s^{-3}<8/d^3$, the full large-s
contribution, divided by E, is below $K_{\rm large}/d$, where
\[
 K_{\rm large}=U_P\frac72\frac{16}{25}\,8
 \sqrt{\frac{v_{\rm high}}{c_{\rm low}P_{\rm lo}}}\,
                  \frac{CB_1}{1+a}I_{\exp}.
\]
The square root uses the corrected smoothing variance factor.

The remaining periodic terms with $n\ge d/2$ are bounded using
$P_d(x)R^d>d^{-5/2}/32$, inverse mode below $2\sqrt d$,
$e^\beta<9/2$, and $E^{-1}<9/2$.
Equation~\eqref{eq:periodic} gives
\[
 1296d^3e^{-d/20}<1/1000\quad(d\ge1001).
\]
This decreases for $d>60$; at 1001 its exact endpoint certificate is
$1296000\cdot1001^3(3/8)^{50}<1$, using $e>8/3$.

The large-saddle verifier certifies, on each whole closed interval,
\begin{equation}\label{eq:finalinterval}
 \frac{K_{\rm small}+K_{\rm pref}+K_{\rm large}}{1001}
       +\frac{306CI_{\exp}}{1001\cdot31}+\frac1{1000}<1.
\end{equation}
Every term bounds all $d\ge1001$, because the algebraic errors decrease in d
and the periodic term has the monotonicity just proved.
Thus this is a uniform dimension proof, not finite dimension sampling.
Together with the split of Section~\ref{sec:split}, it proves
$|[z^m]\Rpoly_d/\kappa-E|<E$ and Proposition~\ref{prop:large}.

\begin{longtable}{c c r}
\caption{The 33 closed intervals and upward-rounded total relative error bounds from \eqref{eq:finalinterval}.}\label{tab:intervals}\\
\toprule Left endpoint & Right endpoint & Certified upper envelope\\\midrule
\endfirsthead
\toprule Left endpoint & Right endpoint & Certified upper envelope\\\midrule
\endhead
\bottomrule\endfoot
$1/10$ & $1/5$ & 0.415204\\
$1/5$ & $3/10$ & 0.354745\\
$3/10$ & $2/5$ & 0.337833\\
$2/5$ & $1/2$ & 0.333091\\
$1/2$ & $3/5$ & 0.333805\\
$3/5$ & $7/10$ & 0.337634\\
$7/10$ & $4/5$ & 0.343535\\
$4/5$ & $9/10$ & 0.349782\\
$9/10$ & $1$ & 0.356797\\
$1$ & $5/4$ & 0.273062\\
$5/4$ & $3/2$ & 0.269584\\
$3/2$ & $7/4$ & 0.268158\\
$7/4$ & $2$ & 0.268035\\
$2$ & $9/4$ & 0.268862\\
$9/4$ & $5/2$ & 0.270358\\
$5/2$ & $11/4$ & 0.272354\\
$11/4$ & $3$ & 0.274745\\
$3$ & $7/2$ & 0.453021\\
$7/2$ & $4$ & 0.450379\\
$4$ & $9/2$ & 0.450842\\
$9/2$ & $5$ & 0.453294\\
$5$ & $11/2$ & 0.457061\\
$11/2$ & $6$ & 0.461875\\
$6$ & $13/2$ & 0.467545\\
$13/2$ & $7$ & 0.473930\\
$7$ & $8$ & 0.911391\\
$8$ & $9$ & 0.904558\\
$9$ & $10$ & 0.903392\\
$10$ & $11$ & 0.906514\\
$11$ & $12$ & 0.912996\\
$12$ & $13$ & 0.922190\\
$13$ & $14$ & 0.933627\\
$14$ & $15$ & 0.946958\\
\end{longtable}


\section{Exact endpoint arithmetic and reproducibility}\label{sec:exact}
This section specifies the analytic tails and enclosure rules used by the
certificates. None of the scalar tests replaces a preceding analytic proof.
For a rational endpoint x,
\[
 q=\left(\frac{1+x}{2}\right)^{-2/(x-1)}\ (x\ne1),\qquad q(1)=e^{-1}.
\]
Exact rational-power bisection encloses q in 48 steps.
The logarithm of a positive rational a is enclosed using
$t=(a-1)/(a+1)$, 128 terms, and remainder
$2|t|^{257}/(257(1-t^2))$.
For exponentials use 97 positive terms (indices 0..96) and the upper remainder
$a^{97}/(97!(1-a/98))$ for $0\le a<20$; negative arguments use reciprocals.
Signed linear combinations in \eqref{eq:vclosed} use the correct endpoint
order. Upward and downward rounding use the dyadic lattice of denominator
$2^{72}$. An upward square-root bound is
\[
 \frac{\lfloor\sqrt{\lfloor a2^{100}\rfloor}\rfloor+1}{2^{50}}.
\]
No floating-point value controls a logical decision.

For the large endpoint sums take $a\le64$ and
$h\le H_a=\max(2,\lfloor128/a\rfloor)$. With $t=q^a$, the four omitted
h tails for $\beta,\beta_1,T_1,T_2$, before multiplication by $v_a$, are
\[
 \frac{t^{H_a}}{(H_a+1)(1-t)},\quad
 \frac{a t^{H_a}}{1-t},\quad
 \frac{a t^{H_a}}{1-t},\quad
 a^2t^{H_a}\left(\frac{H_a}{1-t}+\frac{t}{(1-t)^2}\right).
\]
Let $A>\sqrt{(1+x)/(4\pi)}$ also bound $v_1$ and put $t=q^{65}$.
The omitted a tails are at most
\[
 \frac A{768}\left(1+\frac{t}{2(1-t)}\right),\quad
 \frac A{4(1-t)},\quad
 \frac A{512}\frac{t}{(1-q)(1-t)},\quad
 \frac A8\frac{t}{(1-q)(1-t)^2}.
\]
The first two follow from integral sums after 64; the others use
$a^{-3/2}\le1/512$, $a^{-1/2}\le1/8$ and geometric sums.
For the positive lower log-E sum take $a\le32$ and
$h\le\max(1,\lfloor32/a\rfloor)$, discarding only nonnegative terms.

The three scalar programs use only Python's standard library:
\begin{center}\small\begin{tabular}{l}
\texttt{verify\_small\_x\_majorant.py}\\
\texttt{verify\_majorant\_tail\_constants.py}\\
\texttt{verify\_large\_x\_intervals.py}
\end{tabular}\end{center}
They are preserved byte for byte from the reviewed candidate.
The second imports the third's enclosure helpers. All checks raise explicit
exceptions; they do not use Python \texttt{assert}.
The certificates give respectively the small-saddle scalar/base checks,
four finite-base/induction residuals, and 33 whole-interval rational bounds.
Their exact rational entries and SHA256 digests accompany this paper.
The 6-decimal numbers in Table~\ref{tab:intervals} are rounded
\emph{upwards} from already-certified rational bounds.
The largest exact bound, on $[14,15]$, is
\[
 \frac{2235941074644793697351}{2361183241434822606848}<0.947.
\]
The reproducibility entry point also runs the independent scalar residual
checks, verifies file integrity, and compares the generated witnesses.
Floating-point displays are diagnostics only.

The new proof was manually reviewed and its three exact programs reproduced
independently as part of the accompanying audit. This record is not Lean
verification or external human peer review. The source audit maps every
necessary candidate argument to the final sections; historical incomplete
labels and a corrected scratch variance expression occur only in the
separately identified historical review record.

\section{Completion of the proof}\label{sec:completion}
\begin{proof}[Proof of Theorem~\ref{thm:main}]
For $d\ge1001$, Proposition~\ref{prop:forward} covers
$r\le\lfloor d/4\rfloor$. For the remaining needed coordinates put
$m=d-1-j$, so
\[
 1\le m\le d-2-\lfloor d/4\rfloor<3d/4.
\]
The last strict inequality holds for every residue class of d modulo four.
The elementary estimate
$\mu_d(15)\ge d\nu(15)-15/16>3d/4$ for $d>1000$ follows from
$\log8<209/100$ and
$\nu(15)=15/14-(15/98)\log8>751/1000$.
Thus each required unique saddle lies below 15.
Propositions~\ref{prop:small} and \ref{prop:large} cover all these saddles,
including their common endpoint $1/10$.
Hence $a_{d,j}>0$ for $\lfloor d/4\rfloor+1\le j\le d-2$.
Every remaining tail \eqref{eq:tails} is a positive sum of these coordinates.
The prior endpoint facts and repair criterion give real zeros.
Since the Ehrhart coefficients are nonnegative and $H_d(0)=1$,
$H_d(t)>0$ for $t\ge0$, so all its zeros are strictly negative.

For $d\le1000$, apply de Castro's finite Theorem 1.3.
Its base cases include $H_1=H_2=1$ and $H_3=(1+t)^2$.
This covers all dimensions.
\end{proof}

\subsection{The imported finite certificate}
The finite theorem uses the author's official release
\href{https://github.com/tashimir/ehrhart-132-213-polytopes/releases/tag/v1.0.0}
{\texttt{tashimir/ehrhart-132-213-polytopes}, release \texttt{v1.0.0}},
published on 5 September 2026. The file is
\texttt{refined\_coordinates\_d1\_d1000.csv}, of size 996682200 bytes.
The official release asset digest is
\begin{center}\small\ttfamily
8dbb43cad35659199cd23b5fddc23dd50\\
f76d35adff5b636ac366d460069b5ce
\end{center}
A complete streamed download independently confirmed this byte length and SHA256
digest, and counted 500501 lines including the header. The accompanying
provenance record distinguishes byte-integrity verification from mathematical
reconstruction. We have not independently
reconstructed the CSV or rerun the source's 498500 finite tail inequalities.
Theorem~\ref{thm:main} depends transparently on that prior finite theorem.

\section{A corollary for integral dilates}\label{sec:dilates}
\begin{corollary}\label{cor:dilates}
For every pair of positive integers d,r, the polynomial $h^*(rP_d;t)$ is
constant or has only strictly negative real zeros.
\end{corollary}
\begin{proof}
For $d\ge3$, $P_d$ has dimension $d-1$ in its affine lattice, so the Ehrhart
series has denominator $(1-t)^d$. The Veronese construction gives
\[
 h^*(rP_d;t)=U^d_{r,0}H_d(t),\qquad
 U^d_{r,k}h=\left((1+t+\cdots+t^{r-1})^dh(t)\right)^{\langle r,k\rangle}.
\]
Here $f(t)=\sum_{k=0}^{r-1}t^kf^{\langle r,k\rangle}(t^r)$ defines the
sections. By Theorem~\ref{thm:main}, $H_d$ has nonnegative coefficients and
only negative real zeros. Zhang's Corollary 3.4~\cite[p.~7]{zhang} applies
for every positive integer r, with no degree threshold on r, and makes
the sequence $(U^d_{r,r-1}H_d,\ldots,U^d_{r,0}H_d)$ interlacing.
In particular the zeroth member is real-rooted. Its nonnegative coefficients
and constant term one exclude nonnegative roots.
For $d=1$ the numerator is 1; for $d=2$ the primitive lattice segment
gives $1+(r-1)t$ directly. This is an application of the established Veronese
theorem to this family, not a new preservation theorem.
\end{proof}

\begin{thebibliography}{9}
\bibitem{decastro}
P.~M.~M.~de Castro,
\emph{Ehrhart $h^*$-polynomials of $(132,213)$-avoiding permutation
polytopes: A repair cone and eventual real-rootedness},
arXiv:2609.06096v1, 5 September 2026.
\url{https://arxiv.org/abs/2609.06096v1}.

\bibitem{darroch}
J.~N.~Darroch,
\emph{On the distribution of the number of successes in independent trials},
Ann.\ Math.\ Statist.\ \textbf{35} (1964), 1317--1321, Theorem 4.

\bibitem{pitman}
J.~Pitman,
\emph{Probabilistic bounds on the coefficients of polynomials with only real zeros},
J.\ Combin.\ Theory Ser.\ A \textbf{77} (1997), 279--303.
\url{https://www.stat.berkeley.edu/~pitman/453.pdf}.

\bibitem{zhang}
P.~B.~Zhang,
\emph{Interlacing polynomials and the Veronese construction for rational
formal power series}, author PDF, Corollary 3.4 and equation (1).
\url{https://zhangbiaomath.github.io/papers/2020-prm-ipvc.pdf}.
\end{thebibliography}
\end{document}
